\chapter{Gravité : holonomie, torsion et réponse collective}
\label{chap:gravedad}
\index{gravité}\index{Cartan--Holst}\index{holonomie}
\index{torsion}\index{graviton!non primitif}

La gravité n'est pas introduite ici au moyen d'une formule décimale pour \(G\).
Sa construction réunit deux branches déjà typées. Dans la première, une voie conserve
orientation et mémoire ; une connexion compare des repères en des points distincts ; son
holonomie mesure le retour et sa courbure, le défaut de fermeture. Dans la seconde,
la carte dimensionnelle construit \(G_{\rm int}(\theta_G)\) et
\(\kappa_{\rm int}\) à partir de l'échelle interne d'action, de la vitesse et
de la durée élémentaire. Ce n'est qu'après avoir disposé des deux branches que l'on écrit
l'action Palatini--Cartan--Holst. Aucune valeur métrologique externe ne participe à
cette antériorité.

\section{Retours : position, orientation, phase et mémoire}

Soit \(F_{\rm obs}\) la chronologie observable du transport et
\[
 \mathcal K_{27}=F_{\rm obs}^{27}.
\]
La hiérarchie exacte distingue trois retours :
\[
 \begin{array}{c|c}
 27\ \text{pas}&
 \text{retour de position et inversion d’orientation},\\
 54\ \text{pas}&
 \text{retour de position et d’orientation},\\
 108\ \text{pas}&
 \text{retour complet de la phase observable}.
 \end{array}
\]
Le retour observable n'efface pas le registre accumulé. Si \(g(e)\) compte
une inversion d'orientation et \(s(e)\) un changement effectif de feuillet, alors
\[
 c(\gamma)=\sum_{e\subset\gamma}(g(e),s(e))
\]
satisfait la loi cocyclique de concaténation. Dans les blocs canoniques,
\[
 c_{27}=(1,18),\qquad c_{108}=(4,72).
\]
Par conséquent, la projection observable de \(\mathcal K_{27}\) est d'ordre
quatre, tandis que la coordonnée cumulative continue de croître. Revenir à la
même phase ne signifie pas revenir à la même histoire.

\begin{theorem}[Revêtement cumulatif du retour]
\label{thm:cobertura-retorno-grav-rev6}
Sur \(\mathcal O_{108}\times\mathbb N^2\), le relèvement
\[
 \widetilde F(x,m)=
 \bigl(F_{\rm obs}x,m+c(x\to F_{\rm obs}x)\bigr)
\]
satisfait
\[
 \widetilde{\mathcal K}_{27}(x,m)
 =\bigl(\mathcal K_{27}x,m+(1,18)\bigr)
\]
et
\[
 \widetilde{\mathcal K}_{27}^{\,4}(x,m)
 =\bigl(x,m+(4,72)\bigr)\neq(x,m).
\]
\end{theorem}

\begin{proof}
La constance de \(c_{27}\) sur l'orbite et la loi de concaténation donnent
\[
 c_{108}=\sum_{j=0}^{3}c_{27}(\mathcal K_{27}^jx)
 =4(1,18).
\]
La première coordonnée revient ; la seconde conserve la mémoire accumulée.
\end{proof}

Telle est la donnée qu'une réalisation géométrique doit préserver. On n'identifie pas
la mémoire discrète à la torsion : on construit une représentation
qui permette de les comparer.

\section{Du groupoïde des chemins au groupe de Cartan}

Soit \(\mathcal G_{\rm mem}\) le groupoïde des voies enrichies, avec
composition par concaténation et inverse par renversement. Fixons un élément
\(R_4\in\operatorname{Spin}^+(1,3)\) d'ordre quatre et un vecteur temporel
unitaire \(u\in\mathbb R^{1,3}\) fixé par \(R_4\). Si
\[
c(\gamma)=\bigl(c_g(\gamma),c_s(\gamma)\bigr)\in\mathbb Z^2
\]
est le cocycle additif d'inversion et de changement de feuillet, nous définissons
\[
\rho_{\rm C}^{\rm disc}(\gamma)
=
\left(
R_4^{\,c_g(\gamma)},
\ell_0c_s(\gamma)u
\right)
\in\operatorname{Spin}^+(1,3)\ltimes\mathbb R^{1,3}.
\]

\begin{theorem}[Réalisation discrète du cocycle de retour]
\label{thm:realizacion-discreta-cartan-rev6}
\(\rho_{\rm C}^{\rm disc}\) est un foncteur du groupoïde des voies dans le groupe de
Cartan. En particulier, il conserve l'identité, la concaténation et l'inversion, et
transporte le retour \(27\to54\to108\) sans effacer la coordonnée cumulative.
\end{theorem}

\begin{proof}
La loi du produit semi-direct est
\[
(R,a)(S,b)=(RS,a+Rb).
\]
Comme \(R_4u=u\) et que \(c\) est additif,
\[
\begin{aligned}
\rho_{\rm C}^{\rm disc}(\gamma_2)
\rho_{\rm C}^{\rm disc}(\gamma_1)
&=
\left(
R_4^{c_g(\gamma_2)+c_g(\gamma_1)},
\ell_0\bigl(c_s(\gamma_2)+c_s(\gamma_1)\bigr)u
\right)\\
&=\rho_{\rm C}^{\rm disc}(\gamma_2\gamma_1).
\end{aligned}
\]
La voie identité a un cocycle nul et le renversement change son signe, de sorte
que l'identité et l'inverse sont préservés. Les formules de retour découlent de
\(c_{27}=(1,18)\) et de \(c_{108}=4c_{27}\).
\end{proof}

Ce théorème construit le pont discret. Une réalisation lisse ultérieure
consiste en une représentation
\[
 \rho_{\rm C}:\mathcal G_{\rm mem}
 \longrightarrow\operatorname{Spin}(1,3)\ltimes\mathbb R^{1,3}
\]
et en des champs \(e^I,\omega^{IJ}\) dont l'holonomie entrelace cette représentation.
Le diagramme qui fixe le type est
\[
 \operatorname{Hol}_{\mathcal A_{\ell_0}}
   \bigl(\mathfrak R_{\rm C}(\gamma)\bigr)
 =
 \rho_{\rm C}\bigl(\operatorname{Hol}_{\rm TPK}(\gamma)\bigr).
\]
Les identités, l'inversion, la concaténation et la subdivision doivent être conservées. Cette
équation ne rebaptise pas la retenue courbure : elle exige que le transport
discret déjà représenté et le transport géométrique réalisent la même composition.

Une échelle positive \(\ell_0\) étant fixée, la connexion conjointe est
\[
 \mathcal A_{\ell_0}
 =
 \begin{pmatrix}
 \omega&\ell_0^{-1}e\\
 0&0
 \end{pmatrix}.
\]

\begin{proposition}[Courbure conjointe de Cartan]
\label{prop:curvatura-conjunta-cartan-rev6}
La courbure de \(\mathcal A_{\ell_0}\) est
\[
 \mathcal F_{\ell_0}
 =\dd\mathcal A_{\ell_0}
  +\mathcal A_{\ell_0}\wedge\mathcal A_{\ell_0}
 =
 \begin{pmatrix}
 F_\omega&\ell_0^{-1}T_{\rm tor}\\
 0&0
 \end{pmatrix},
\]
où
\[
 F_\omega^{IJ}
 =\dd\omega^{IJ}+\omega^I{}_K\wedge\omega^{KJ},
 \qquad
 T_{\rm tor}^{I}=D_\omega e^I.
\]
\end{proposition}

\begin{proof}
On multiplie les matrices de formes en utilisant le produit extérieur. Le
bloc diagonal produit \(F_\omega\) ; le bloc de translation produit
\(\dd e+\omega\wedge e=D_\omega e\).
\end{proof}

La courbure et la torsion sont deux composantes d'un même défaut de transport,
mais remplissent des fonctions distinctes. La courbure mesure le retour des repères ;
la torsion mesure la fermeture de la cotétrade.

\section{Gauss--Stokes et capacité surfacique}

Dans un complexe cellulaire fini, pour toute \(1\)-cochaîne \(\Omega\) et toute
\(2\)-chaîne \(z\),
\[
 \boxed{\quad
 \langle\delta\Omega,z\rangle
 =\langle\Omega,\partial z\rangle.
 \quad}
\]
L'identité est la définition du cobord comme adjoint du bord.
Lorsque \(z\) somme les faces orientées de manière cohérente d'une
pseudovariété, les arêtes intérieures s'annulent et seul subsiste le
bord géométrique. La même égalité a déjà été lue comme
vorticité--circulation et comme courbure--holonomie.

L'échelle surfacique possède également un modèle exact. Dans le graphe cubique
\(N\times N\times N\), toute coupe entre deux faces opposées contient au moins
\(N^2\) arêtes et une couche transversale atteint la borne. Pour \(N=9\),
\[
 9^3=729=27^2.
\]
Cette égalité cardinale relie volume et coupe minimale ; elle n'identifie pas les
groupes \((\mathbb Z/9)^3\) et \((\mathbb Z/27)^2\), qui ont des exposants
distincts.

\section{L'action Palatini--Cartan--Holst}

Soit \(M\) une variété lisse orientée de dimension quatre, munie d'une cotétrade non
dégénérée \(e^I\), d'une connexion de Lorentz
\(\omega^{IJ}=-\omega^{JI}\), et de conditions au bord annulant le terme
variationnel de bord. Écrivons
\[
 \Sigma^{IJ}=e^I\wedge e^J,\qquad
 (\star X)^{IJ}=\frac12\epsilon^{IJ}{}_{KL}X^{KL},
 \qquad
 G_{\rm int}=G_{\rm int}(\theta_G),\qquad
 \kappa_{\rm int}=\frac{8\pi G_{\rm int}}{c_{\rm int}^{4}}.
\]
Ces deux dernières quantités ne sont pas définies au moyen de l'action : elles proviennent de
la construction constitutive de la \cref{eq:G-interno}. Le
\cref{thm:homogeneidad-pch} démontre auparavant que la combinaison suivante
appartient à la droite d'action.
L'action du premier ordre est
\[
\boxed{
 S[e,\omega,\Psi]
 =\frac1{2\kappa_{\rm int}c_{\rm int}}\int_M
 \Sigma_{IJ}\wedge
 \left(\star F^{IJ}+\frac1\gamma F^{IJ}\right)
 -\frac{\Lambda_{\rm int}}{\kappa_{\rm int}c_{\rm int}}
  \int_M\operatorname{vol}_e
 +S_{\rm m}[e,\omega,\Psi].}
\]
Le paramètre interne construit au chapitre précédent spécialise
\[
 \gamma=\Gamma_{6\to8}.
\]
La substitution est effectuée après avoir obtenu
\(\Gamma_{6\to8}\) à partir de \(A,C^\ast,q_\pm\) et de l'inclusion
hexade--octade ; l'action n'est pas utilisée pour redéfinir ces entrées.

\begin{theorem}[Équation de Cartan--Holst]
\label{thm:cartan-holst-rev6}
Si le courant de spin est normalisé au moyen de
\[
 \delta_\omega S_{\rm m}
 =-\frac12\int_M\delta\omega^{IJ}\wedge\sigma_{IJ},
\]
la variation par rapport à \(\omega\) produit
\[
 \boxed{\quad
 D_\omega
 \left[(\star+\gamma^{-1}I)\Sigma^{IJ}\right]
 =\kappa_{\rm int}c_{\rm int}\,\sigma^{IJ}.
 \quad}
\]
\end{theorem}

\begin{proof}
On utilise \(\delta F^{IJ}=D_\omega\delta\omega^{IJ}\), on intègre par parties et
on annule le coefficient de la variation arbitraire. Le terme
cosmologique ne dépend pas de \(\omega\).
\end{proof}

En signature lorentzienne, \(\star^2=-I\) sur les bivecteurs. Pour
\(\gamma\in\mathbb R\setminus\{0\}\),
\[
 (\star+\gamma^{-1}I)^{-1}
 =\frac{\gamma^2}{1+\gamma^2}
  (-\star+\gamma^{-1}I).
\]
Par conséquent, si le courant de spin s'annule,
\[
 T_{\rm tor}^I=0,\qquad
 \omega=\mathring\omega(e).
\]
En présence de matière spinorielle, la torsion est déterminée algébriquement par le
courant. Elle ne constitue pas, dans l'action minimale, un degré propagatif
indépendant.

\section{Équation métrique, spin et bilan géométrique}

Pour fixer sans ambiguïté la normalisation, nous définissons les sources mesurées dans
la droite d'action par
\[
 \delta_g S_{\rm m}
 =-\frac12\int_M\operatorname{vol}_e\,
 \mathcal T^{S}_{\mu\nu}\,\delta g^{\mu\nu}.
\]
Lorsque les coordonnées de la tétrade ont la dimension d'une longueur, nous écrivons
\[
 T^{\rm phys}_{\mu\nu}:=c_{\rm int}\mathcal T^{S}_{\mu\nu},
 \qquad
 U^{\rm phys,spin}_{\mu\nu}
 :=c_{\rm int}\mathcal U^{S,\rm spin}_{\mu\nu}.
\]
Après élimination algébrique de la connexion, l'équation métrique adopte les
deux formes équivalentes
\[
 \boxed{
 G_{\mu\nu}+\Lambda_{\rm int}g_{\mu\nu}
 =\kappa_{\rm int}c_{\rm int}
 \bigl(\mathcal T_{\mu\nu}^{S,\rm LC}
       +\mathcal U_{\mu\nu}^{S,\rm spin}\bigr)
 =\kappa_{\rm int}
 \bigl(T_{\mu\nu}^{\rm phys,LC}
       +U_{\mu\nu}^{\rm phys,spin}\bigr).}
\]
La contribution de spin est quadratique en le courant correspondant et
dépend du couplage de matière ; son coefficient n'est pas fixé sans déclarer
l'action fermionique et les conventions de dualité.

Les identités géométriques
\[
 D_\omega T_{\rm tor}^{I}=F^{I}{}_{J}\wedge e^{J},
 \qquad
 D_\omega F^{IJ}=0
\]
impliquent, après élimination de la torsion,
\[
 \nabla^\mu\bigl(T_{\mu\nu}^{\rm phys,LC}
             +U_{\mu\nu}^{\rm phys,spin}\bigr)=0.
\]
Ce bilan de Bianchi--Noether est une condition de cohérence du
transducteur gravitationnel : aucun terme supplémentaire ne peut être incorporé sans
provenir d'une action compatible ou satisfaire le même bilan.

\section{Correspondance surfacique–volumétrique \(\pi/\varphi^2\)}

Le calendrier \((\Sigma,\Pi,\Pi)\) produit la matrice d'incidence
\[
 F_{\rm av}=\begin{pmatrix}0&1\\1&1\end{pmatrix}.
\]
Sa mesure de Parry donne
\(\mu_{\rm ar}/\mu_{\rm vol}=\varphi^{-2}\), et le lecteur régional marque une
seule fois la fermeture \(\pi_{\rm HMT}\). Le retour marqué satisfait
\[
 \Theta_n
 =\pi_{\rm HMT}\frac{F_{n-1}}{F_{n+1}}
 \longrightarrow
 \frac{\pi_{\rm HMT}}{\varphi^2}.
\]
Ce rapport n'est pas une autre version de \(G\). C'est la correspondance adimensionnelle entre
lecture surfacique et lecture volumétrique. Sous une même amplitude positive
\(\mathcal M\),
\[
 \mathcal I_{\rm av}=\mu_{\rm vol}\mathcal M,\qquad
 \mathcal G_{\rm av}=\pi_{\rm HMT}\mu_{\rm ar}\mathcal M
\]
donnent exactement
\[
 \frac{\mathcal G_{\rm av}}{\mathcal I_{\rm av}}
 =\frac{\pi_{\rm HMT}}{\varphi^2}.
\]
Sur le feuillet plan, la restriction diagonale est
\(\mathcal G_{\rm tor}=\mathcal I_{\rm tor}\). L'équivalence locale et
l'ambivalence de bord sont donc deux régimes du lecteur, non deux
valeurs incompatibles d'une masse.

\section{Typage de \(G\) comme transducteur gravitationnel et ordre causal de ses lecteurs}

Le type et l'homogénéité de \(G\) appartiennent à la construction
constitutive préalable. La \cref{eq:G-interno} fixe une seule fois la famille
covariante \(G_{\rm int}(\theta_G)\), et le
\cref{thm:homogeneidad-pch} démontre que son insertion dans l'action possède
la dimension d'une action. Ce chapitre ne redéfinit pas cette famille : il utilise sa
section avec l'holonomie déjà construite. La causalité n'est pas linéaire,
mais convergente :
\[
\begin{aligned}
 \mathcal B_{\rm geom}&:
 \text{retour et cocycle}\to\text{représentation de Cartan}
 \to(e,\omega),\\
 \mathcal B_{\rm dim}&:
 (\hbar_{\rm int},c_{\rm int},t_0,\theta_G)
 \to G_{\rm int}\to\kappa_{\rm int},
\end{aligned}
\]
\[
 (\mathcal B_{\rm geom},\mathcal B_{\rm dim})
 \to S_{\rm PCH}
 \to\text{torsion algébrique, équation métrique et bilan}.
\]
La comparaison métrologique intervient à la fin et ne sélectionne aucune des
deux branches.

\Needspace{22\baselineskip}
\section{Origine collective du mode gravitationnel dans l'holonomie des voies}

L'action précédente a pour objets fondamentaux la cotétrade et la connexion.
La torsion de Cartan est algébrique ; la courbure possède le secteur propagatif qui
contient les ondes gravitationnelles. Le formalisme n'a donc pas besoin de
postuler un graviton élémentaire pour définir l'interaction. Une éventuelle
quantification des perturbations peut produire un mode collectif
d'hélicités \(\pm2\) ; ce mode serait une excitation du transport
géométrique, non une donnée ajoutée au niveau d'APP, de TRIT ou de TPK.

Cette distinction exclut également une confusion fréquente. L'existence
d'ondes classiques démontre une propagation de courbure ; elle ne démontre pas à elle seule
un quantum fondamental. Réciproquement, décrire le champ au moyen de la
géométrie n'élimine pas ses degrés radiatifs. HMT--MD situe la question au
bon endroit : on construit d'abord la mémoire et son holonomie ; on étudie ensuite
le spectre de ses excitations.
